\section{Guidance Algorithm and Results}\label{sec:results}
\subsection{Algorithm}\label{sec:algorithm}
The guidance law has three steps. 1) Once, before the descent, the time of flight is set to the smallest $t_f$ at which
Eqs.~\eqref{eq:prog} are feasible, found by bisection in about ten solves, plus a margin of 0.25~s. 2) At every guidance step
the program is solved for the remaining time of flight from the current state estimate; nothing is carried over from the
previous plan, and no initial trajectory, trust region, or penalty weight is required. 3) At every control step the reference
position, velocity, attitude, body rate, throttle, and gimbal angle are evaluated from Eqs.~\eqref{eq:flat} in closed form
and passed to the tracking controller. If a solve is infeasible or its plan fails the slew-rate check, the flight time is
extended by up to 1~s, and otherwise the previous plan, which remains valid, is kept. For $S=16$ the program has 63 spline
coefficients; as passed to the solver it has about 1000 variables and 480 second-order cones, of the types used by 3-DOF
LCvx, and one solve takes about 40~ms.

\subsection{Setup}
All results are numerical and use the vehicle of Table~\ref{tab:vehicle}. The baseline, denoted S, is a standard 6-DOF SCvx
implementation \cite{mao2016,szmuk2018} on the quaternion rigid body with 35 intervals, virtual control, and a trust region,
with the same planning limits. The program of Eqs.~\eqref{eq:prog} is denoted F. Both use the Clarabel interior-point solver
on a four-core desktop processor. The simulated plant is the quaternion rigid body with mass-dependent inertia and lever arm,
a $\pm15$~deg gimbal box with a 100~deg/s slew limit, a 400 to 1000~N throttle with a 2000~N/s slew limit, wind of up to
3~m/s with gusts, 2\% thrust noise, a 5\% error in specific impulse, and gimbal misalignment. One tracking controller, a
linear-quadratic regulator in center-of-percussion coordinates at 50~Hz, serves both planners. The test set is 40 in-air
descents (altitude 40 to 60~m, offset up to 21~m, initial tilt up to 3~deg) and 30 hops launched from the pad at 5 to 6~m/s
toward a point up to 4~m away. A landing succeeds if the touchdown speed is at most 1~m/s, the lateral error at most 1~m,
and the tilt at most 5~deg.

\subsection{Single Plans: Cost, Exactness, and Soundness}
On eight in-air scenarios at equal final time, F uses 0.10\% more fuel than S on average and 0.28\% at most, with one SOCP
and a median solver time of 0.046~s against 25 to 35 SOCPs and 1.30~s. Its minimum feasible final time is 0.04 to 0.22~s
longer, and up to 0.41~s longer on 4-s hops; this is the incompleteness of Proposition~\ref{prop:inc}. To test Step~1 of the
proof, the throttle and gimbal commands of F were applied without feedback to the quaternion rigid body with depleting mass
and constant $\varepsilon$, integrated to a relative tolerance of $10^{-13}$: in 16 replays the vehicle came to rest within
23~nm of the target, which is integration error. To test Theorem~\ref{thm:sound}, 300 random problems were each solved with
the constants above and with randomly chosen ones; 458 of the 600 solves were feasible, and no plan violated a planning limit
at any of 1201 instants. To test Proposition~\ref{prop:inc}, both planners and a nonconvex solver were run on 20 hops of
exactly 4~s. No flyable plan was found for 11 of the targets, and F is infeasible at all 11; it is also infeasible at five of
the nine targets where a plan exists. A 3-DOF LCvx program reports an optimal solution at all 20, which is the false
feasibility that Theorem~\ref{thm:A} predicts for a complete program.

\subsection{Closed Loop with Computation Time Counted}
Simulation studies of guidance commonly apply each plan at the instant its state was sampled. Here a plan computed from the
state at time $t$ becomes active at $t$ plus its solver time, and each planner replans as soon as it can: S at most every
0.5~s and F every 0.1~s. Table~\ref{tab:closed} and Fig.~\ref{fig:td} give the results. Both planners land every case. F is more accurate
laterally by more than an order of magnitude, touches down more slowly, and moves the gimbal a third as far. On hops it uses
5\% more fuel, because it flies them slightly longer. With every solver time multiplied by four, to represent a slower flight
computer, S fails four of the 30 hops, and F, replanning every 0.2~s, fails none. The source of the advantage should be
stated carefully: when both planners replan at 1~Hz, plans are applied instantly, and both fly the same final time, the
lateral accuracy is the same (medians of 3.1~cm in the air and 2.4~cm on hops for each). The accuracy in
Table~\ref{tab:closed} comes from the replanning rate, which a solver time of 30~ms permits and a solver time of 1~s does not.

%%CLOSEDTABLE%%
%%TDFIG%%
\subsection{Limitations}
F is worse than S in fuel on short hops and in minimum flight time. The slew-rate limits are outside
Theorem~\ref{thm:sound}; fewer than 1\% of replans failed the check and were not flown. Theorem~\ref{thm:sound} is proved here
and tested numerically but is not machine-checked, apart from Eq.~\eqref{eq:cop} for constant mass, and it has not had an
independent reader as Theorem~\ref{thm:Bstar} has. The assumption of constant $\varepsilon$ must be confirmed from measured mass
properties. Nothing here has been run on flight hardware, with actuator lag, or with state-estimation error, and the SCvx
baseline was not retuned for speed. On the impossibility side, the proved gap is small for descents that begin in the air, and
for the gimbaled rigid body the result is certified only at four final times under simplifying assumptions.

\section{Conclusion}
For powered descent with a finite-rate attitude, a single convex program of the kind lossless convexification belongs to
cannot be exact: an exact program would make the set of landable initial states convex, and a dive to one side followed by a
brake, together with its mirror image, shows that the set is not convex at any final time. This is machine-checked. What a
single convex program can still be is sound. For a lander steered by one gimbaled engine, writing the problem for the center
of percussion removes the coupling between steering and translation and leaves the actuator limits as the only nonconvex
part; restricting those limits, instead of relaxing them, yields one second-order cone program whose plans are flyable at
every instant by proof. It costs a tenth of a percent of fuel and a few tenths of a second of minimum flight time, has no
initial guess, trust region, or iteration count, and solves fast enough to replan at the rate of the control system, which in
our simulations mattered more for touchdown accuracy than the last fraction of a percent of optimality. Future work is a
certificate of nonconvexity that accounts for the gimbal-rate limit and depleting mass, an independent check and
formalization of the soundness proof, slew-rate limits inside the guarantee, and flight.

\section*{Acknowledgments}
% AIAA policy (Ethical Standards, Sec. 3) requires a description of AI use here and a disclosure at submission.
% The first two sentences below are a factual description of how this work was produced. Complete the rest truthfully.
An artificial-intelligence system (Claude, Anthropic) was used extensively in this work. It derived and wrote the proofs,
wrote the Lean formalization and the planning and simulation software, ran the numerical experiments, and drafted the text
and figures of this manuscript. [State here what the Contributing Student Authors did themselves, for example: posed the
question, defined the vehicle data, directed the investigation, and checked the derivations, the code, and the results.] The
independent check of the impossibility proof was carried out by [name of the reviewer, with permission]. [State any other
contributors and sponsors.] The authors are responsible for the content of the paper.
